\section{Saponification de l'éthanoate d'éthyle}
\textbf{Source du sujet~:} Baccalauréat 2022.

\subsection{Introduction}
On souhaite déterminer le temps de demi-réaction de la saponification de l'éthanoate d'éthyle par deux méthodes~: un suivi cinétique pH-métrique (partie A) puis l'exploitation mathématique de la loi de vitesse (partie B). L'équation de la réaction modélisant cette transformation chimique est~:

\small
\[
\begin{array}{ccccccc}
\raisebox{-0.5\height}{% pubchem 8857
\chemfig{H_3C-[:30]C(=[:90]\charge{60=\|,120=\|}{O})-[:330]\charge{90=\|,-90=\|}{O}-[:30]CH_2-[:330]CH_3}} &
\raisebox{-0.5\height}{+} &
\raisebox{-0.5\height}{\chemfig{H-\charge{90=\|,0=\|,270=\|}{O}^{\ominus}}} &
\raisebox{-0.5\height}{$\xrightarrow{}$} &
\raisebox{-0.5\height}{% pubchem 175
\chemfig{\mcfright{\charge{30=\|,210=\|,300=\|}{O}}{^{\ominus}}-[:120]C(=[:60]\charge{30=\|,90=\|}{O})-[:180]{CH_3}}} &
\raisebox{-0.5\height}{+} &
\raisebox{-0.5\height}{% pubchem 702
\chemfig{OH-[:210,,1]{CH_2}-[:150]{CH_3}}} \\
\multicolumn{1}{c}{\text{Éthanoate d'éthyle}} & & \multicolumn{1}{c}{\text{Ion hydroxyde}} & & \multicolumn{1}{c}{\text{Ion éthanoate}} & & \multicolumn{1}{c}{\text{Éthanol}}
\end{array}
\]
\normalsize

\subsection{Données}
\begin{itemize}
    \item Masse volumique de l'éthanoate d'éthyle~: $\rho_{\text{é.é.}} = 0,925~\mathrm{g/mL}$.
    \item Masse molaire de l'éthanoate d'éthyle~: $M_{\text{é.é.}} = 88,0~\mathrm{g/mol}$.
    \item Constante d'équilibre de l'autoprotolyse de l'eau~: $K_{\mathrm{e}} = 10^{-14}$.
\end{itemize}

\subsection{Partie A~: Suivi pH-métrique de la saponification}

\begin{wrapfigure}{r}{0.4\textwidth}
    \centering
    \includegraphics[width=\linewidth]{ste7_ethanoate_ph.png}
%    \caption{Suivi pH-métrique de la saponification de l'éthanoate d'éthyle}
    \AMClabel{fig:ethanoate_ph}
\end{wrapfigure}

Le suivi est réalisé en mettant en œuvre le protocole suivant~:

\begin{itemize}
    \item Dans un bécher, verser~:
    \begin{itemize}
        \item 10,0 mL de soude à 0,10 mol/L~;
        \item environ 35 mL d'eau~;
        \item environ 15 mL d'éthanol~;
        \item 2,0 mL d'éthanoate d'éthyle.
    \end{itemize}
    \item Déclencher l'acquisition informatique de la mesure du pH toutes les secondes. Le déclenchement définit l'instant initial $t = 0~\mathrm{s}$.
\end{itemize}

\WFclear

\vspace{0.5cm}
\begin{questionmult}{saponifA1}
    Déterminer les valeurs des quantités de matière initiales d'éthanoate d'éthyle et de soude.\\
    \begin{EnvQuadrillage}[NbCarreaux=19x6,Grille=Seyes,Marge=1]
    \end{EnvQuadrillage}
    \AMCOpen{}{
        \mauvaise[F]{NA}\scoring{b=0}
        \bonne{Idée}
        \bonne{Éthanoate id}
        \bonne{Éthanoate}
        \bonne{Soude id}
        \bonne{Soude}
    }
    \explain{
        Pour déterminer les quantités de matière initiales d'éthanoate d'éthyle et de soude, on procède comme suit :
        \textbf{Pour la soude} :
        La concentration molaire \(C\) est définie par la relation \(C = \frac{n}{V}\), où \(n\) est la quantité de matière (en \(\mathrm{mol}\)) et \(V\) le volume de la solution (en \(\mathrm{L}\)).
        On peut donc exprimer la quantité de matière \(n\) comme suit :
        \[
        n_{\ce{HO^-}} = C \times V.
        \]
        En utilisant les valeurs données, \(C = \SI{0,10}{\mol\per\litre}\) et \(V = \SI{0,010}{\litre}\), on obtient :
        \[
        n_{\ce{HO^-}} = \SI{0,10}{\mol\per\litre} \times \SI{0,010}{\litre} = \SI{1,0e-3}{\mol}.
        \]
        \textbf{Pour l'éthanoate d'éthyle} :
        La quantité de matière \(n\) est déterminée à partir de la masse \(m\) et de la masse molaire \(M\) selon la relation \(n = \frac{m}{M}\).\\
        La masse \(m\) est calculée à partir de la masse volumique \(\rho\) et du volume \(V\) :
        \[
        m = \rho \times V = \SI{0,925}{\gram\per\milli\litre} \times \SI{2,0}{\milli\litre} = \SI{1,85}{\gram}.
        \]
        En utilisant la masse molaire \(M = \SI{88,0}{\gram\per\mol}\), on obtient la quantité de matière :
        \[
        n = \frac{m}{M} = \frac{\SI{1,85}{\gram}}{\SI{88,0}{\gram\per\mol}} = \SI{2,10e-2}{\mol}.
        \]
    }
\end{questionmult}

Dans ces conditions, on considère que la concentration en éthanoate d’éthyle reste
constante et égale à sa valeur initiale tout au long de la transformation chimique. La loi de vitesse de réaction peut alors s’écrire comme celle d’une loi d’ordre~1~:  $v = k_1 \times [\ce{HO-}]$ dans laquelle $k_1$ est appelée constante de vitesse apparente.

On souhaite établir la relation entre le pH mesuré et la concentration en quantité de matière $C$ en ions hydroxyde dans le mélange réactionnel au même instant.

\vspace{0.5cm}
\begin{questionmult}{saponifA2}
    Écrire l'équation de la réaction d'autoprotolyse de l'eau.\\
    \begin{EnvQuadrillage}[NbCarreaux=19x1,Grille=Seyes,Marge=1]
    \end{EnvQuadrillage}
    \AMCOpen{}{
        \mauvaise{NA}\scoring{b=0}
        \bonne{ID}
        \bonne{OK}
    }
    \explain{
    L'équation de l'autoprotolyse de l'eau est~:
    \ce{2 H2O <=> H3O^+ + HO^-}
    }
\end{questionmult}

\vspace{0.5cm}

\begin{questionmult}{saponifA3}
    Donner l'expression de la constante d'équilibre \(K_{\mathrm{e}}\) de la réaction d'autoprotolyse de l'eau en fonction des concentrations à l'équilibre.\\
    \begin{EnvQuadrillage}[NbCarreaux=19x1,Grille=Seyes,Marge=1]
    \end{EnvQuadrillage}
    \AMCOpen{}{
        \mauvaise{NA}\scoring{b=0}
        \bonne{ID}
        \bonne{OK}
    }
    \explain{
    La constante d'équilibre \(K_{\mathrm{e}}\) est donnée par~:
    $$
    K_{\mathrm{e}} = [\ce{H3O^+}] \times [\ce{HO^-}]
    $$
    }
\end{questionmult}

\begin{questionmult}{saponifA4}
    En déduire la relation entre la concentration $C$ en ions hydroxyde et le pH mesuré à tout instant.\\
    \begin{EnvQuadrillage}[NbCarreaux=19x6,Grille=Seyes,Marge=1]
    \end{EnvQuadrillage}
    \AMCOpen{}{
        \mauvaise{NA}\scoring{b=0}
        \bonne{ID}
        \bonne{OK}
    }
    \explain{
    À partir de la constante d'équilibre de l'eau $K_e = [\ce{H3O^+}] \times [\ce{HO^-}]$, et en utilisant la définition du pH ($(\text{pH} = -\log [\ce{H3O^+}]$), on peut exprimer la concentration en ions hydroxyde comme suit~:
    $$
    [\ce{HO^-}] = \frac{K_e}{[\ce{H3O^+}]}
    $$
    En utilisant la relation \(\text{pH} = -\log [\ce{H3O^+}]\), on obtient~:
    $$
    [\ce{HO^-}] = 10^{\text{pH} - 14}
    $$
    }
\end{questionmult}

À partir des valeurs du pH mesuré à chaque seconde et en utilisant la relation entre C et le
pH, on obtient le graphique Figure~\AMCref{fig:ethanoate_graph}.

\begin{figure}[h]
    \centering
    \includegraphicscorr[width=\linewidth, keepaspectratio]{ste7_ethanoate_graph}
    \caption{Graphique de la concentration en ions hydroxyde en fonction du temps}
    \AMClabel{fig:ethanoate_graph}
\end{figure}

\vspace{0.5cm}
\begin{questionmult}{saponifA5}
     Déterminer, à partir de la figure~\AMCref{fig:ethanoate_graph}, le temps de demi-réaction $t_{1/2}$ , exprimé en seconde, de la saponification. Une justification graphique est attendue.\\
    \begin{EnvQuadrillage}[NbCarreaux=19x4,Grille=Seyes,Marge=1]
    \end{EnvQuadrillage}
    \AMCOpen{}{
        \mauvaise{NA}\scoring{b=0}
        \bonne{ID}
        \bonne{OK}
    }
    \explain{
    Pour déterminer le temps de demi-réaction \(t_{1/2}\) graphiquement~:
    \begin{itemize}
        \item Repèrer la concentration $C_0$ sur le graphique. Ici, $C_0 = 0,016$ mol/L.
        \item Repérer le point où $C(t) = \frac{C_0}{2} = \SI{0,008}{\mol\per\liter}$.
        \item Lire la valeur de \(t\) correspondant à ce point sur l'axe des abscisses~: ici $t_{1/2} = \SI{40}{\second}$.
    \end{itemize}    
    }
\end{questionmult}

\subsection{Partie B~: Mathématiques}

On note $C(t)$ la concentration en ions hydroxyde, exprimée en mol/L, à l’instant $t$, exprimé en seconde, et $C_0$ la concentration en ions hydroxyde à l’instant $t = 0$.\\
Dans les conditions décrites dans la partie A, $C_0 = 0,016$ mol/L et $k_1 = 0,017$ s$^{-1}$.\\
La fonction $C$ est donc solution de l’équation différentielle (E) suivante~:
$$ y' = -k_1 y \quad \text{(E)} $$

\begin{questionmult}{saponifB1}
    Vérifier que la fonction $C$ définie sur $[0 ; +\infty[$ par $C(t) = C_0 e^{-k_1 t}$ est une solution de (E).\\
    Montrer que $C(0) = C_0$. On admet que $C$ est la seule solution de (E) qui vérifie $C(0) = C_0$.
    \begin{EnvQuadrillage}[NbCarreaux=19x6,Grille=Seyes,Marge=1]
    \end{EnvQuadrillage}
    \AMCOpen{}{
        \mauvaise{NA}\scoring{b=0}
        \bonne{verifID}
        \bonne{verifOK}
        \bonne{c0ID}
        \bonne{c0OK}
    }
    \explain{
    Pour vérifier que $C(t) = C_0 e^{-k_1 t}$ est une solution de (E), dérivons $C(t)$~:
    $$
    C'(t) = -k_1 C_0 e^{-k_1 t} = -k_1 C(t)
    $$
    Donc, $C(t)$ satisfait bien l'équation différentielle (E). De plus, à $t=0$~:
    $$
    C(0) = C_0 e^{-k_1 \times 0} = C_0
    $$
    }
\end{questionmult}

\begin{questionmult}{saponifB2}
    Déterminer par le calcul le temps de demi-réaction $t_{1/2}$. On donnera la valeur exacte, puis l’arrondi à la seconde. Interpréter ce résultat dans le contexte de l’exercice.
    \begin{EnvQuadrillage}[NbCarreaux=19x6,Grille=Seyes,Marge=1]
    \end{EnvQuadrillage}
    \AMCOpen{}{
        \mauvaise{NA}\scoring{b=0}
        \bonne{ID}
        \bonne{Interpr}
        \bonne{Interpr}
    }
    \explain{
    Le temps de demi-réaction $t_{1/2}$ est défini comme le temps nécessaire pour que la concentration en ions hydroxyde soit divisée par deux. On a donc~:
    $$
    C(t_{1/2}) = \frac{C_0}{2} = C_0 e^{-k_1 t_{1/2}}
    $$
    En prenant le logarithme, on obtient~:
    $$
    t_{1/2} = \frac{\ln(2)}{k_1} = \frac{\ln(2)}{0,017} \approx 40,8 \text{ secondes}
    $$
    Ce résultat signifie qu'il faut environ 41 secondes pour que la concentration en ions hydroxyde soit réduite de moitié. Ce qui correspond à la valeur trouvée sur le graphique figure~\AMCref{fig:ethanoate_graph}.
    }
\end{questionmult}
